\scanpage{42}
Soit $S$ un schéma de base, et soit $X$ un schéma de Br. Severi sur $S$ muni d'une \uline{trisection} non ramifiée $X'$, i.e un sous-schéma \inserted{fermé} de $X$ qui est non ramifié sur $X$, d'ordre 3 sur $X$.

Un tel truc peut être considéré comme un objet à schéma structural $\mathfrak S_3$, opérant sur $P^1_S$ comme il a été dit. La classification de tels objets est donc la $=$ que la classification des \deleted{\unsure{rev}} $\mathrm H^1(S,\mathfrak S_3)$ des \unsure{revêtem\textsuperscript{ts}|revêt\textsuperscript{ts}|rev.} \inserted{gal} \uline{non ramifiés} de \deleted{\unsure{schéma}} groupe $\mathfrak S_3$ de $S$. De plus, \unsure{on|$\sigma$} définit dans $X$ un $S$-automorphisme \inserted{involutif}, noté $\sigma$, qui transforme $X'$ en un $X''$, dit \uline{associé} à \unsure{$X$!|$X'$.|$X'$!} \unsure{Évidemt|Évidemment}, si on \unsure{partait|partant} de $(X,X'')$, on aurait le $=$ $\sigma$, et on \unsure{réobtient|restitue|trisection} $X'$ [à préciser --- ]. De plus, on obtient [par descente --- ] une \uline{bisection}, non ramifiée \deleted{[\unsure{du moins}} \unsure{en les pts|sur les pts} de $S$ de car. $\ne 2,3$]..... On peut passer au quotient, \deleted{\unsure{qui}} par $\mathfrak S_3$, \deleted{\unsure{pour}} on trouve \unsure{que|pour} :
\[
P^1_S\xrightarrow{p}P^1_S
\]
(\unsure{comp. au|lorsque} $\mathfrak S_3$ \unsure{opérant|opère}
% archive sheet number: 41
